\begin{python0}
from solutions import *; clear()
\end{python0}
\sectionthree{Why arrays?}

The above explains only the syntax and how to work with arrays. It also explains that if you need to work with a lot of variables, performing the same operation on these variables (such as printing their values), then you might want to use an array.

But \ldots probably the most important ``view'' of an array is that you should view it as a \EMPHASIZE{storage container} of values. Why is that important?

You see containers of values everywhere. For instance when you buy something online (say at amazon.com), amazon has to retrieve information on the product. The product information is in a container -- the products catalog database. At this point we only have a container (or array) of integer values or double values or boolean values or characters. Much later you will learn about containers of not just single values but aggregate values such as a container of customer data made up of firstname, lastname, email address, cellphone number, etc.

An array allows us to store computation so that a program can save time when the program runs and avoid recomputation all the time.

Read the above several times!!!

So how does an array ``store computation'' and ``avoid recomputation''?

Remember that we already know how to check if a number is prime: A positive integer n greater than 1 is a prime if we test n \% d for d = 2, 3, 4, \ldots, n -- 1. If n \% d is 0, then d is a divisor and therefore n is not a prime. If n \% d is not zero for all d = 2, 3, 4, \ldots, n -- 1, then n must be a prime.

The next thing I told you is that you can improve on the speed of the above algorithm by testing only d = 2, 3, 4, \ldots up to including the square root of n. For the case of n = 101, you only need to test d = 2, 3, 4, \ldots, 10 instead of d = 2, 3, 4, \ldots, 100. That's a saving of 90\% \ldots !!!

But \ldots there's \texttt{yet} another way to speed up the program: you only need to test n \% d for d running through \texttt{primes} up to the square root of n. Therefore for n = 101, you only need to test n \% d for d = 2, 3, 5, 7. So from d = 2, 3, 4, \ldots, 10 you go down to d = 2, 3, 5, 7, from 9 values of d to 4. That's about 50\% \ldots !

How would an array help?

One way is to compute an array of 1000 primes before you run your program. Let's say the array is called:

\begin{center}
\texttt{int prime[1000];}
\end{center}

After that for every n you want to test for primality, you simply test n \% prime[i] for i = 0, 1, 2, \ldots and stop when prime[i] if greater than the square root of n. The only caveat is that this array is finite. So if you need to check for n \% d when d is larger than the $1000^{th}$ prime then you need to perform some tests without using the \texttt{prime} array. For instance, since the $1000^{th}$ prime is 7919, and the square root of 1357931 is 1165.30, to check if 1357931 is a prime we test

\tab[3em]{1357931 \% 2}\\
\tab[3em]{1357931 \% 3}\\
\tab[3em]{1357931 \% 5}\\
\tab[3em]{\vdots}\\
\tab[3em]{1357931 \% 1637}

(the prime after 1637 is 1657 which is greater than 1165.30) and of course the primes 2, 3, 5, \ldots, 1637 are taken from our \texttt{prime} array. By the way 1637 is the $259^{th}$ prime. Therefore there are 259 checks:

\tab[3em]{1357931 \% 2}\\
\tab[3em]{1357931 \% 3}\\
\tab[3em]{1357931 \% 5}\\
\tab[3em]{\vdots}\\
\tab[3em]{1357931 \% 1637}\\

Using our crude method of testing

\tab[3em]{1357931 \% 2}\\
\tab[3em]{1357931 \% 3}\\
\tab[3em]{1357931 \% 4}\\
\tab[3em]{1357931 \% 5}\\
\tab[3em]{\vdots}\\
\tab[3em]{1357931 \% 1357930}\\

required 1357929 checks. If we test up to the square root of 1357929:

\tab[3em]{1357931 \% 2}\\
\tab[3em]{1357931 \% 3}\\
\tab[3em]{1357931 \% 4}\\
\tab[3em]{1357931 \% 5}\\
\tab[3em]{...}\\
\tab[3em]{1357931 \% 1165}\\

we have to perform 1164 checks. In summary:

\begin{itemize}
\item
  Test 1357931 \% d for d = 2, 3, 4, \ldots, 1357930: 1357929 tests
\item
  Test 1357931 \% d for d = 2, 3, 4, \ldots, 1165: 1164 tests
\item
  Test 1357931 \% d for d = 2, 3, 5, \ldots, 11637: 259 tests
\end{itemize}

See the difference?

The square of 7919 is 62710561. Therefore with an array of 1000 primes, we can check if n is a prime for n up to 62710561. If n is greater than 62710561, then we need a larger \texttt{prime} array.

This is also the reason why Google works so fast. When you search for web pages containing ``dog'', Google would send you the relevant URLs in a split second because Google continually scans for billions of web pages, precomputes, and stores keywords of the web pages before you even ask for ``dog'' related web pages. The only difference is that Google stores the results in databases and not arrays.

